% ============================================================================
%   APPENDIX D — CODE AND DATA AVAILABILITY
% ============================================================================

\chapter{Code and Data Availability}
\label{app:code}

This appendix documents the archival of the code, prompts, prediction 
files, and analysis outputs produced during this study, in accordance 
with the reproducibility commitments described in 
Section~\ref{sec:met-reproducibility}.

% ----------------------------------------------------------------------------
\section{Repository}
\label{sec:app-repo}

The complete project is archived in a public GitHub repository:

\begin{center}
    \url{https://github.com/RegisLikassi/llm-clinical-trial-calibration}
\end{center}

The repository is organized into four top-level directories, each 
serving a distinct purpose:

\begin{itemize}
    \item \texttt{thesis/} contains the \LaTeX{} sources of this 
    manuscript, including the main file, all chapters, appendices, 
    and the bibliography.
    \item \texttt{code/} contains the Kaggle notebook used for the 
    twelve model runs and the Python scripts that perform the 
    downstream analyses (calibration audit, hypothesis tests, 
    recalibration, conformal prediction, mode collapse analysis, and 
    composite utility scoring).
    \item \texttt{data/} contains the raw prediction files (twelve 
    Parquet files, one per configuration) and the CSV outputs of 
    every downstream analysis reported in Chapter~\ref{chap:results}.
    \item \texttt{docs/} contains the monthly progress reports 
    submitted during the project, the defense slides, and this 
    appendix in standalone form.
\end{itemize}

% ----------------------------------------------------------------------------
\section{Data Availability}
\label{sec:app-data}

The primary benchmark used in this study, 
TrialGPT-Criterion-Annotations, is publicly available under a CC BY 
4.0 license as a supplementary release of Jin et 
al.~\cite{jin2024trialgpt}. It can be downloaded directly from the 
TrialGPT repository:

\begin{center}
    \url{https://github.com/ncbi-nlp/TrialGPT}
\end{center}

The twelve prediction files produced during this study are stored in 
the \texttt{data/predictions/} subdirectory of the project 
repository, one Parquet file per configuration. Each file contains, 
for the $1{,}015$ patient--criterion pairs, the following columns: 
the annotation identifier, the criterion type, the predicted label, 
the associated confidence, the expert consensus label, the GPT-4 
label, and the full softmax distribution over the six candidate 
labels. All downstream analysis outputs --- calibration metrics, 
permutation-test results, recalibration comparisons, conformal 
prediction sets, mode collapse statistics, and composite utility 
scores --- are archived as CSV files in the \texttt{data/analyses/} 
subdirectory.

% ----------------------------------------------------------------------------
\section{Reproducing the Experiments}
\label{sec:app-reproduce}

Independent reproduction of the pipeline follows five steps:

\begin{enumerate}
    \item Clone the repository and install the Python dependencies 
    listed in \texttt{code/requirements.txt}.
    \item Download the TrialGPT-Criterion-Annotations dataset from 
    the URL above and place it in the \texttt{data/raw/} directory.
    \item Obtain access to the six selected models from the 
    HuggingFace Hub. Llama-3.1 requires acceptance of Meta's 
    community license, obtained through the HuggingFace interface.
    \item Execute the Kaggle notebook (or its local equivalent) to 
    produce the twelve prediction files. Random seeds are fixed at 
    $42$ throughout the pipeline.
    \item Execute the analysis scripts in the order documented in 
    the notebook's \texttt{README.md} to reproduce every table and 
    figure of Chapter~\ref{chap:results}.
\end{enumerate}

% ----------------------------------------------------------------------------
\section{Computational Requirements}
\label{sec:app-compute}

The twelve model runs require a GPU with at least $16$~GB of memory 
(a Tesla~T4 or equivalent) and approximately ten to fifteen hours of 
GPU time. The downstream analyses run on CPU and complete in a few 
minutes on standard hardware. All experiments reported in this 
thesis were performed on the Kaggle free-tier compute environment, 
which is publicly accessible and imposes no barriers to independent 
replication.

% ----------------------------------------------------------------------------
\section{Software Versions}
\label{sec:app-versions}

The pipeline was tested with the following software versions:

\begin{itemize}
    \item Python 3.10
    \item PyTorch 2.1 with CUDA 12.1
    \item \texttt{transformers} 4.44
    \item \texttt{bitsandbytes} 0.43 (for 4-bit NF4 quantization)
    \item \texttt{scikit-learn} 1.5
    \item \texttt{numpy} 1.26, \texttt{pandas} 2.2, \texttt{scipy} 1.13
    \item \texttt{matplotlib} 3.9 (for figure generation)
\end{itemize}

The exact version pinning is available in 
\texttt{code/requirements.txt} at the head of the repository. The 
commit hashes of the notebook and analysis scripts at the time of 
submission are recorded in \texttt{docs/submission\_commits.txt}.